\section{文献综述}
\label{sec:lit}

多智能体与强化学习组合管理。把组合决策表述为连续控制问题的研究由深度强化学习
框架系统化\cite{jiang2017}，此后 FinRL 等开源环境把市场模拟、算法与回测流程整合为
统一实验平台\cite{liu2020}，多策略协调研究进一步在收益策略与风险策略之间引入
动态协调机制\cite{li2024}。算法层面的基础工作确立了确定性策略梯度与自举式价值估计的
可行性\cite{mnih2015,lillicrap2016,sutton2018}，金融场景的系统评述与端到端组合优化实现
可分别参见 \cite{charpentier2023} 与 \cite{cong2021}。这类研究的共同设定是：
以历史数据训练一个状态到权重的映射，并用样本外表现作为评价标准。
其目标函数通常是绝对财富或绝对收益，在满仓长仓约束下该目标存在结构性退化——
与决策无关的基准收益项主导梯度，剩余梯度主要来自换手惩罚，机制收敛于静态等权。
本文分离出这一退化并给出主动收益的修正方案，使自适应机制具备识别互补性的条件
（命题 \ref{prop:degeneracy}）。

状态自适应资产配置。Hamilton 的马尔可夫状态切换模型为刻画收益分布的制度变化
提供了基础工具\cite{hamilton1989}，Guidolin 和 Timmermann 将其用于资产配置，
发现状态切换改善样本外表现\cite{guidolin2007}。这一传统隐含「状态可识别即有配置价值」
的判断。本文证明该判断不充分：即使状态被完美识别，若信号有效性不随状态变化
（$\Gamma=0$），门控的价值仍为零（推论 \ref{cor:identification}）。
状态识别的价值取决于互补性与持续性两个条件，本文将两者写成可估计的对象，
为「何时使用状态切换模型」提供了可计算的判据。

预测组合与「组合之谜」。预测组合文献记录了一个稳健事实：
简单平均预测经常打败样本内最优的估计权重组合\cite{timmermann2006}；
相关经典结果包括协方差矩阵收缩\cite{ledoit2004}与等权组合的样本外稳健性
\cite{demiguel2009}。既有解释多归于「估计误差」，但缺少可计算的临界值。
本文给出的阈值不等式补上这一环：$\Theta$ 就是「组合之谜」的临界线，
$\hat\Gamma<\Theta$ 时估计最优权重不如简单平均。渠道分解进一步说明，
二阶矩比一阶矩更容易估计，以风险配置为主的合成本质上比以收益择时为主的合成
更容易获得稳健收益。这为「简单策略为何难以被超越」提供了一个新的解释维度。

互补性的第一个来源：信息成本与注意力配置。门控价值来自互补性，
而互补性需要经济学来源。Grossman 和 Stiglitz 指出，若价格已充分反映信息，
收集信息便无利可图，市场因此不可能完全有效\cite{grossman1980}；
Sims 的理性疏忽理论说明，决策者处理信息的能力受有限信道容量约束，
会把注意力配置给成本低、精度高的信息源\cite{sims2003}；
DellaVigna 和 Pollet 的证据显示注意力不足会延迟信息进入价格\cite{dellavigna2009}。
这条线索对本文有两个推论。其一，同一信号在不同市场状态下的信息含量不同：
状态决定了哪些公开信息与收益更相关，信号的预测精度随状态变化，互补性因此为正。
其二，状态信息本身有价格：估计条件信息系数需要按状态切分样本，
每格的估计精度随样本量下降，这一估计误差成本正是本文阈值中的 $\Theta$。

互补性的第二个来源：行为偏差的状态依赖。Odean 记录的处置效应使投资者
过早卖出盈利、过久持有亏损\cite{odean1998}；Barber 和 Odean 记录过度自信
提高交易频率并降低净收益\cite{barberodean2001}；Hong 和 Stein 的统一理论说明，
信息在不同投资者群体间的逐步扩散可以同时产生动量和反转\cite{hongstein1999}。
这类偏差的强度随市场状态变化，上涨行情中的过度自信与下跌行情中的处置效应更为突出，
以动量和反转为基础构造的信号，其有效性因此带有状态依赖。
本文的专家信号集合正由这类价量信号构成，互补性因此不是设定假设，而是可估计的对象。

渠道差异的理论对应：有限套利。「一阶互补性有限、二阶互补性为主」
这一发现的对应理论是有限套利与信息的估计成本差异。基于收益水平（一阶矩）的套利
需要承担基本面风险与噪声交易者风险，在资金约束下难以充分实施\cite{shleifer1997}；
而波动率与相关性（二阶矩）可以从公开价格序列以较低成本估计。
理性疏忽的预测正是把注意力配置给后者，本文的渠道分解为该预测提供了
横截面资产配置场景下的证据。

机器学习资产定价与稳健性检验。机器学习在收益预测中的应用已成体系。
Gu 等在美国股票横截面上比较十三类模型，报告树集成与神经网络的样本外 $R^2$
相对线性模型提升逾一倍\cite{gu2020}，并在自编码器与深度学习方法上做了拓展
\cite{gu2021,chen2024}。与之并行的是对多重检验的批评：Harvey 等要求新因子满足
更高的 $t$ 值门槛\cite{harvey2016}，Feng 等在给定已有因子的条件下检验新增因子的
定价信息\cite{feng2020}，Kozak 等指出高维因子的定价信息由少数主成分承载
\cite{kozak2020}，Kelly 等讨论模型复杂度与样本量的匹配关系\cite{kelly2024,kelly2023}。
中国市场的机器学习量化投资与文本情绪信号已有系统证据\cite{libin2019,jiang2021,shen2019}，
金融机器学习的方法论综述则强调回测过拟合、信息泄漏与前视偏差的风险
\cite{lopezdeprado2018}。本文属于该传统的稳健性检验分支，检验对象是模型组合机制本身：
本文要求报告「自适应优势」的研究同时给出主动暴露、互补性与噪声底线的对比、
以及渠道分别检验。三组诊断量的计算成本低，作用是把机制奏效与机制退化区分开。

人工智能在金融系统中的治理含义。机器学习在信贷、支付与投顾环节的渗透
使可解释性与公平性成为监管议题。Fuster 等发现算法定价在降低成本的同时
对部分群体产生差异化影响\cite{fuster2022}；基础模型的能力与风险综述指出，
模型规模扩张不自动带来可靠性\cite{bommasani2021}；中央银行视角的研究关注
机器学习在监管科技中的应用与边界\cite{chakraborty2017,bis2024aer}。
本文的阈值不等式与三组诊断量回应这一议题：它把「人工智能组合是否真的更优」
转化为可在同一数据集上独立复算的问题，为算法审计提供可执行的判据。

表 \ref{tab:lit-map} 归纳本文与各组文献的差异边际。

\begin{table}[htbp]
\centering
\caption{文献群、缺口与本文增量}
\label{tab:lit-map}
\begin{threeparttable}
\footnotesize
\begin{tabular}{p{2.5cm}p{5.4cm}p{5.3cm}}
\toprule
文献群 & 已有贡献与留下的缺口 & 本文增量 \\
\midrule
多智能体与强化学习组合管理 & 建立了状态到权重的映射框架；
目标函数多为绝对收益，未识别退化解 & 分离出满仓约束下的退化机制，
给出主动收益修正（命题 \ref{prop:degeneracy}） \\
状态自适应资产配置 & 证明状态切换可改善配置；
默认「状态可识别即有配置价值」 & 给出门控价值的充要条件，
说明识别状态不等于有价值（推论 \ref{cor:identification}） \\
预测组合与组合之谜 & 记录简单平均的稳健优势，
以估计误差笼统解释 & 给出临界值的闭式解与比较静态，
把经验规律变为可复算判据（定理 \ref{thm:threshold}） \\
信息经济学与行为金融 & 说明信息的成本、注意力的稀缺与偏差的状态依赖；
未落到多信号合成的决策 & 把互补性追溯到信息成本与行为偏差，
并以渠道分解检验其可识别性（定理 \ref{thm:channels}） \\
机器学习资产定价与稳健性 & 建立了预测与多重检验的规范；
检验对象是单一模型 & 把稳健性检验对象扩展到组合机制，
提出三个低成本诊断量 \\
人工智能金融治理 & 关注可解释性、公平性与监管科技的边界；
缺少可操作的判定程序 & 提供可在同一数据集复算的阈值判据，
把审计要求落为具体计算 \\
\bottomrule
\end{tabular}
\end{threeparttable}
\end{table}
